%%%%%%%%%%%%%%%%%%%%%%%%%%%%%%%%%%%%%%%%%%%%%%%%%%%%%%%%%%%%%%%%%%%%%%%%%%%%%%%%
%%  EINSTELLUNGEN – die einzige Datei, in der du deine Angaben einträgst.
%%  Wird von ALLEN Hauptdateien (Endabgabe + Zwischenabgaben) eingelesen.
%%  Entspricht dem Block "Die Angaben" der offiziellen Gerüste
%%  pp-portfolio-geruest-2026-2027.tex / pp-portfolio-geruest-einlage-2026-2027.tex
%%%%%%%%%%%%%%%%%%%%%%%%%%%%%%%%%%%%%%%%%%%%%%%%%%%%%%%%%%%%%%%%%%%%%%%%%%%%%%%%

%%%%%%%%%%%%%%%%%%%%%%%%%%%%%%%%%%%%% Die Angaben %%%%%%%%%%%%%%%%%%%%%%%%%%%%%%%
\newcommand{\Semester}{WH 2026-2027}
\newcommand{\Projekt}{Dynamik Keltischer Wackelsteine}   % Titel laut Themenliste WS 26/27
%% Schwerpunkt abgeleitet aus der Aufgabenverteilung im Arbeitsplan (1.4 Bewegungsgleichungen,
%% 1.6 Simulation Umkehrdynamik, 2.1.3/2.1.4 Firmware + Kalibrierung, 5.1 Datenpipeline) – prüfen!
\newcommand{\Schwerpunkt}{Bewegungsgleichungen und numerische Simulation der Spin-Umkehr}
\newcommand{\SchwerpunktKurz}{Simulation der Spin-Umkehr}        % kurze Form für die Kopfzeile
\newcommand{\IhrName}{Kristijan Manevski}                         % Vorname Nachname
\newcommand{\Matrikelnummer}{\platzhalter{Matrikelnummer}}
\newcommand{\Studiengang}{\platzhalter{Studiengang + Semester}}   % z.B. Physik B.Sc., 5. Semester
\newcommand{\emailAdresse}{\platzhalter{Uni-E-Mail-Adresse (\dots @uni-ulm.de)}}
%% Abgabedatum: jede Zwischenabgabe-Hauptdatei setzt ihr eigenes \dasDatum (Moodle-Frist);
%% ohne Vorgabe gilt das Datum der Endabgabe.
\providecommand{\dasDatum}{15.\,03.\,2027}

%% Weitere Angaben, die in mehreren Einlagen gebraucht werden
\newcommand{\Partner}{Janosch Lustig}                             % für Abgrenzung der Schwerpunkte
%% laut Titelseite von Janoschs Arbeitsplan-Abgabe (28.10.2026)
\newcommand{\PartnerSchwerpunkt}{Umkehrzeit und Parameterabhängigkeit}
\newcommand{\Laborbuchlink}{https://link-zum-elektronischen-laborjournal}  % Pflicht in Einlage 3!
%% Einlage 3: ein Praktikumstag (Do 10--12 Uhr, O26/233 bzw. O26/234), ≠ Tag von Janosch.
%% Termine: 15.10., 22.10., 29.10., 05.11., 12.11., 19.11., 26.11., 03.12., 10.12., 17.12.2026,
%%          14.01., 21.01., 28.01., 04.02., 11.02.2027 (Abgabe Protokoll: 09.12.2026)
\newcommand{\Protokolldatum}{TT.MM.2026}
\newcommand{\PaperKey}{KondoNakanishi2017}      % Einlage 2: Kondo & Nakanishi, Phys. Rev. E 95, 062207
\newcommand{\Vortragsdatum}{28.\,01.\,2027}      % Einlage 5
\newcommand{\FremdVortragKey}{VORLAGE-Vortrag}  % Einlage 6: bib-Key des anderen Vortrags

%%%%%%%%%%%%%%%%%%%%%%%%%%%%%%%%%%%%%%%%%%%%%%%%%%%%%%%%%%%%%%%%%%%%%%%%%%%%%%%%
%%  MODUS-SCHALTER
%%
%%  \Abgabefalse  = Entwurf: grüne Hilfen sichtbar, Kontrollseite mit
%%                  Umfangsmessung, Platzhalter rot markiert.
%%  \Abgabetrue   = Abgabe: grüne Farbe -> schwarz (wie offizielle Zeile 15/16),
%%                  alle Hilfen + Kontrollseite + Umfangsboxen verschwinden.
%%                  Übrig gebliebene \platzhalter{} bleiben ROT sichtbar und
%%                  erzeugen eine Warnung -> dann ist noch etwas offen!
%%
%%  \Hinweisefalse blendet im Entwurf nur die grünen Hilfstexte aus
%%                  (nützlich, um den Umfang realistisch zu messen).
%%%%%%%%%%%%%%%%%%%%%%%%%%%%%%%%%%%%%%%%%%%%%%%%%%%%%%%%%%%%%%%%%%%%%%%%%%%%%%%%
\newif\ifAbgabe
\newif\ifHinweise

\Abgabefalse      % <-- für die Abgabe auf \Abgabetrue setzen
\Hinweisetrue     % <-- Hilfstexte im Entwurf ein-/ausblenden
